\documentclass{article}

\usepackage{amsmath,amssymb}
\usepackage{xurl}
\usepackage[hidelinks]{hyperref}
\setlength{\emergencystretch}{2em}

% Shared commands for notes in docs/paper-gaps/.

\DeclareUnicodeCharacter{2080}{\ifmmode{}_{0}\else\textsubscript{0}\fi}
\DeclareUnicodeCharacter{2081}{\ifmmode{}_{1}\else\textsubscript{1}\fi}
\DeclareUnicodeCharacter{2082}{\ifmmode{}_{2}\else\textsubscript{2}\fi}
\DeclareUnicodeCharacter{2099}{\ifmmode{}_{n}\else\textsubscript{n}\fi}

\newcommand{\Lean}{\textsc{Lean}}
\ExplSyntaxOn
\NewDocumentCommand{\leanid}{m}{
  \ifmmode
    \textsf{\detokenize{#1}}
  \else
    \group_begin:
    \urlstyle{sf}
    \tl_set:Nn \l_tmpa_tl {#1}
    \tl_replace_all:Nnn \l_tmpa_tl {\_} {_}
    \exp_args:NV \path \l_tmpa_tl
    \group_end:
  \fi
}
\ExplSyntaxOff
\providecommand{\tr}{\operatorname{tr}}
\newcommand{\tnleanIssueBase}{https://github.com/LionSR/TNLean/issues/}
\newcommand{\tnleanPrBase}{https://github.com/LionSR/TNLean/pull/}
\newcommand{\tnleanDocBase}{https://sirui-lu.com/TNLean/docs/}
\newcommand{\ghissue}[1]{\href{\tnleanIssueBase#1}{GitHub issue \##1}}
\newcommand{\ghpr}[1]{\href{\tnleanPrBase#1}{PR \##1}}
\newcommand{\doclink}[1]{\href{\tnleanDocBase#1.html}{\path{#1}}}

% Dirac notation and the identity operator, as in blueprint/src/macros/common.tex.
% \providecommand keeps note-local definitions authoritative.
\usepackage{dsfont}
\providecommand{\ket}[1]{|#1\rangle}
\providecommand{\bra}[1]{\langle #1|}
\providecommand{\braket}[2]{\langle #1 | #2 \rangle}
\providecommand{\ketbra}[2]{|#1\rangle\!\langle #2|}
\providecommand{\Id}{\mathds{1}}
\providecommand{\one}{\mathds{1}}
\providecommand{\C}{\mathbb{C}}
\providecommand{\MN}[1]{M_{#1}(\C)}
\providecommand{\MD}{M_D(\C)}

% Verdict marker: \gapnote{<kind>}{<status>}. Kind is the policy
% classification (clarification, local-correction, scope-restriction,
% unfaithful, false-source, open-gap); status is open, wip (elimination
% underway), resolved, or historical. Machine-read by the site index;
% prints a verdict line.
\newcommand{\gapnote}[2]{\par\noindent\textbf{Verdict:} #1 (#2).\par}


% Fallbacks keep this note compilable when a different command.tex is found first.
\providecommand{\leanid}[1]{\textsf{\detokenize{#1}}}
\providecommand{\gapnote}[2]{\par\noindent\textbf{Verdict:} #1 (#2).\par}

\title{The Constant Bond Dimension in the Algorithm of Landau, Vazirani and Vidick}
\author{}
\date{2026-10-25}

\begin{document}
\maketitle
\gapnote{local-correction}{resolved}

\section{The source statement}

Landau, Vazirani and Vidick~\cite{Landau2015Polynomial} consider
$H=\sum_{i=1}^{n-1}H_i$ on $n$ qudits of dimension $d$, with $H_i$ acting on
the sites $i,i+1$ and $0\le H_i\le 1$, ground state $\Gamma$ and gap
$\varepsilon$ (arXiv:1307.5143v1, \path{algo53.tex}, line~283). Quoting
Hastings~\cite{Hastings2007AreaLaw} and Arad, Kitaev, Landau and
Vazirani~\cite{Arad2013AreaLaw}, they state (lines~342--345) that for every
cut $(i,i+1)$ and every $\delta$ there is a constant
\begin{align}
    B_\delta=\exp\Bigl(O\bigl(\tfrac1\varepsilon\log^3d\,\log\tfrac1\delta\bigr)\Bigr)
    \label{eq:lvv13_constant_bond_source}
\end{align}
such that the normalized trimmed vector
$v=\operatorname{trim}^i_{B_\delta}\Gamma/
    \lVert\operatorname{trim}^i_{B_\delta}\Gamma\rVert$ satisfies
$\lvert\langle\Gamma,v\rangle\rvert\ge1-\delta$. No proof is given in the
source.

\section{What is formalized}

The formalized statement concerns a chain Hamiltonian $H=\sum_Xh_X$ on $L$
sites of dimension $q\ge1$, with Hermitian terms supported on sets of
diameter at most $R$ and $\lVert h_X\rVert\le J$, and a gapped ground state
$\Omega$ with gap $\Delta>0$. The source Hamiltonian is the case $R=1$,
$J=1$. Let $C$ be the bound on the entropy of an interval given by the
one-dimensional area law of the library, let $b$ be its cut budget constant,
which depends on $q$ and $R$ only, and let
$s=\sqrt{(1+J/\Delta)\,b}$. For $\delta>0$ put
\begin{align}
    B_\delta=\Bigl\lfloor\exp\Bigl(C+32\,s\log\frac{2e^{e/2}}{\delta}\Bigr)
    \Bigr\rfloor .
    \label{eq:lvv13_constant_bond_formal}
\end{align}
Then for every cut $k$ and every choice of the kept Schmidt vectors, the
normalized trimmed vector $v$ satisfies
\begin{align}
    \operatorname{Re}\langle v,\Omega\rangle\ge1-\delta ;
    \label{eq:lvv13_constant_bond_overlap}
\end{align}
the overlap is real and nonnegative.\footnote{\leanid{TNLean.MPS.AreaLaw.ChainHamiltonian.one_sub_le_re_inner_chainTrim}
    in \doclink{TNLean/MPS/AreaLaw/ViableSet/ConstantRankCut}.}

The proof is short. The marginal tail bound of the area law gives, at every
cut and every $w\ge0$, a spectral projection of the left marginal $\rho$ of
rank at most $e^{C+w}$ discarding mass at most $2e^{e/2}e^{-w/(32s)}$. The
choice $w=32\,s\log(2e^{e/2}/\delta)$ makes this mass equal to $\delta$ and
the rank at most $B_\delta$. By Ky Fan's maximum principle the sum of the
$B_\delta$ largest eigenvalues of $\rho$ is at least $1-\delta$; this sum is
$\lVert\operatorname{trim}\Omega\rVert^2$, and the overlap in
(\ref{eq:lvv13_constant_bond_overlap}) is $\lVert\operatorname{trim}\Omega
    \rVert\ge\sqrt{1-\delta}\ge1-\delta$.

\section{The constant}

The bound (\ref{eq:lvv13_constant_bond_formal}) has the structure of
(\ref{eq:lvv13_constant_bond_source}): it depends on neither the length of
the chain nor the cut, and it is the exponential of an affine function of
$\log(1/\delta)$, rounded down. Its constants are different. The slope is $32s$ and the
intercept involves $C$, both taken from the area law proved in the library,
not the dependence $\varepsilon^{-1}\log^3d$ of
(\ref{eq:lvv13_constant_bond_source}). The later steps of the source use
only that $B_\delta$ is independent of $n$ and of the cut, so the
different constant affects the explicit running time, not the conclusions.
The theorems hold under the source's hypotheses: its nearest-neighbour
Hamiltonian with $0\le H_i\le1$ is a chain Hamiltonian of range $1$ and
strength $1$.

\section{The range of \texorpdfstring{$\delta$}{delta}}

The source states Lemma \path{constantbondapprox} and Corollary
\path{c:1} for any constant $\delta$, and both are formalized for every
$\delta>0$. In the lemma, for $\delta\ge1$ the conclusion holds because the
overlap is nonnegative. No upper bound on $\delta$ is needed in the
corollary, because the approximation can be chosen at a smaller error.

For Corollary \path{c:1} of the source (lines~361--364), let $\delta>0$
and apply the constant bond dimension at $\delta_0=\min\{\delta,1/12\}$.
Writing $E_0$ for the ground-state energy, the corrected trimming-energy
estimate gives
\begin{align}
    \langle v,Hv\rangle
    \le E_0+\frac{4\sqrt{\delta_0}}{1-4\delta_0}
    \le E_0+6\sqrt{\delta_0}\le E_0+6\sqrt\delta.
    \label{eq:lvv13_constant_bond_corollary_energy}
\end{align}
The overlap is at least $1-\delta_0\ge1-\delta$, and the bond dimension
$B_{\delta_0}$ remains independent of the chain length and cut. Thus the
corollary retains its stated energy and overlap bounds for every
$\delta>0$, despite the correction needed in Lemma \path{l:2}.

\section{Nontrivial chains in the viable leading step}

The later viable-witness and leading-vector statements currently assume
$L\ge1$ and $q\ge2$. These restrictions are separate from the constant-bond
approximation above. They allow the estimate $\Delta\le4$ for a range-one
Hamiltonian with interaction norm at most one, which in turn controls
$c=(\Delta/169)^2$ and the energy error in the leading-vector argument.
For $q=1$, the chain Hilbert space is one-dimensional and the state and
energy conclusions can instead be proved directly, without a gap upper
bound. A complete extension must also specify a bond cutoff of at least
one. The empty chain requires a separate formulation of the quantities
$c/L$ and $c/(2L)$. Proving these separate statements is the plan for
extending the later construction; it does not supply the source's
quantitative AKLV bound. This remaining extension is part of the viable-set
construction tracked in issue~\href{https://github.com/LionSR/TNLean/issues/9820}{9820}.

\section{The source's constant}

The explicit dependence $\varepsilon^{-1}\log^3d$ is not formalized. It would
follow from the bound of Arad, Kitaev, Landau and Vazirani,
arXiv:1301.1162~\cite{Arad2013AreaLaw}, which gives the Schmidt rank in
(\ref{eq:lvv13_constant_bond_source}), by the argument of the previous
sections.

\bibliographystyle{alpha}
\bibliography{references}

\end{document}
